\documentclass{article}

\PassOptionsToPackage{numbers,sort&compress}{natbib}
\usepackage[dblblindworkshop]{neurips_2026}
\workshoptitle{Sim2Science}

\usepackage[utf8]{inputenc}
\usepackage[T1]{fontenc}
\usepackage{hyperref}
\usepackage{url}
\usepackage{booktabs}
\usepackage{array}
\usepackage{amsmath}
\usepackage{amssymb}
\usepackage{graphicx}
\usepackage{microtype}
\usepackage{xcolor}
\usepackage{yhmath}

\newcommand{\vect}[1]{\boldsymbol{#1}}
\newcommand{\Ldata}{\mathcal{L}_{\mathrm{data}}}
\newcommand{\Ldesign}{\mathcal{L}_{\mathrm{design}}}
\newcommand{\todo}[1]{\textcolor{red}{[TODO: #1]}}
\newcommand{\method}{DiffOED}
\newcolumntype{L}[1]{>{\raggedright\arraybackslash}p{#1}}

\title{Differentiable Physics for Sensor Placement in Full-Waveform Inversion}

% The dblblindworkshop option suppresses this block in the review PDF.
\author{
  Arnaud Mercier\textsuperscript{1} \quad
  Ben Moseley\textsuperscript{2} \quad
  Henrik Rasmus Thomsen\textsuperscript{1} \\
  \textsuperscript{1}Institute of Geophysics, ETH Zurich, Zurich, Switzerland \\
  \textsuperscript{2}Department of Earth Science and Engineering, Imperial College London, London, UK \\
  \texttt{\{arnaud.mercier,henrik.thomsen\}@eaps.ethz.ch} \\
  \texttt{b.moseley@imperial.ac.uk}
}

\begin{document}

\maketitle

% \begin{abstract}
% Optimal experimental design (OED) for full-waveform inversion (FWI) often selects sources or receivers from a predefined candidate survey using a criterion derived from a local linearization. This restricts the design to a preset grid and may misalign the design criterion with the finite nonlinear inversion that it is intended to improve. We instead optimize receiver coordinates continuously by differentiating through a finite sequence of time domain FWI updates. The resulting online design loop alternates between short blocks of FWI, receiver-position updates, and reacquisition. We compare three design objectives. Oracle model error and waveform misfit differentiate through the unrolled FWI trajectory, whereas stochastic Gauss--Newton illumination differentiates only through wave solver sensitivities. We evaluate the methods on a two dimensional acoustic synthetic example inspired by guided-wave non-destructive testing. Across ten paired multi-seed repetitions, all three optimized layouts reduce model error relative to the fixed layout baseline. Mean reductions range from 9.4 to 18.6\% under the same budget of 100 FWI updates. These results provide preliminary evidence that differentiable OED can adapt receiver locations to a finite inversion procedure without requiring a predefined candidate grid.
% \end{abstract}

% \begin{abstract}

% Full-waveform inversion (FWI) uses wave simulators to infer material properties
% from measured waveforms, but its reconstruction quality and convergence depends strongly on
% where sensors are placed. Conventional optimal experimental design often
% selects sensors positions from a predefined candidate survey using a criterion derived
% from a local linearization. This restricts the design to a preset grid and may misalign the design criterion with the finite nonlinear inversion that it is intended to improve. To tackle this issue, we introduce a bilevel active acquisition framework that
% optimizes continuous receiver coordinates using differentiable framework.
% Each outer iteration assimilates observations from a real world measurement device through a finite block of FWI
% updates, updates the receiver positions, and reacquires data at the new
% geometry.

% We investigate two objectives that do not require knowledge of the
% true material model: a waveform-misfit objective differentiated through the
% unrolled FWI trajectory and a stochastic Gauss--Newton illumination objective
% differentiated only through wave-solver sensitivities. We evaluate the framework
% on a two-dimensional acoustic synthetic example inspired by guided-wave
% non-destructive testing. Across ten paired repetitions and an equal budget of
% 100 FWI updates, waveform-misfit and illumination-based placement reduce mean
% model error by 9.4\% and 16.7\%, respectively, relative to a fixed receiver
% layout. Although this proof of concept uses synthetic observations and does not
% yet test simulator misspecification, its acquisition loop is designed to replace
% synthetic reacquisition with physical measurements. The framework therefore
% provides a path for using an imperfect differentiable simulator to guide where
% new real-world data should be collected.

% \end{abstract}

\begin{abstract}


% In full-waveform inversion (FWI), the sensor positions that best
% constrain the material model are generally unknown in advance, yet poor sensor
% placement can substantially degrade reconstruction quality. In physical settings, this challenge is compounded by discrepancies between the wave simulator used for inversion and the observed system. Rather than fixing sensor positions in advance using the simulator alone, we present an differentiable active acquisition loop that alternates between short blocks of FWI, sensor-position updates, and reacquisition. The loop couples differentiable wave simulation with adaptive data acquisition, allowing the sensor geometry to evolve with the reconstruction instead of being selected once from a predefined candidate set. 


In full-waveform inversion (FWI), reconstruction quality depends strongly on sensor placement, yet informative sensor positions are generally unknown in advance. This challenge is particularly acute in physical settings, where the wave simulator used for inversion may differ from the observed system. We present a differentiable active-acquisition loop that alternates short blocks of FWI with continuous receiver-coordinate updates and reacquisition, allowing the geometry to adapt to the evolving reconstruction without requiring a predefined candidate set. We investigate two acquisition design objectives that do not require the true
material model: waveform misfit differentiated through the unrolled FWI
trajectory and stochastic Gauss--Newton illumination differentiated only through
wave-solver sensitivities. In a two-dimensional synthetic acoustic example
inspired by guided-wave non-destructive testing, both objectives improve
reconstruction relative to a fixed receiver layout across ten paired
repetitions. Under the same budget of 100 FWI updates, the waveform misfit and
stochastic illumination objectives reduce mean model error by 9.4\% and 16.7\%,
compared to the fixed layout. These synthetic results motivate future closed-loop experiments in which physical measurements replace simulated reacquisition.

\end{abstract}


\section{Introduction}

Full-waveform inversion reconstructs material properties by repeatedly solving a wave equation and minimizing waveform mismatch~\citep{virieux2009overview}. Its reconstruction quality depends strongly on where sources and receivers are placed, yet measurements and wave simulations are expensive. Optimal experimental design (OED) therefore seeks a small, informative acquisition geometry. In seismic FWI, established methods typically compute sensitivities for a comprehensive candidate survey and greedily select measurements according to spectral or resolution criteria derived from the Gauss--Newton system~\citep{maurer2017oed,krampe2021oed,mercier2025combined}. These approaches are effective, but they require a discrete candidate set, often use greedy selection strategies that can lead to suboptimal designs, and optimize a proxy rather than the actual nonlinear reconstruction outcome.

Differentiable simulators allow a different formulation: treat acquisition
coordinates as continuous variables and optimize them jointly with the FWI workflow. Bilevel learning has been used to select
regularization parameters in inverse problems~\citep{holler2018bilevel} and to
place space--time observations in variational data
assimilation~\citep{castro2020observation}. Related differentiable experimental
design methods use physics-informed or neural
reconstructors~\citep{hemachandra2025pied,darges2025node,liu2025dspo}. Most
closely, \citet{downing2026bilevel} optimizes sensor coordinates for
frequency-domain FWI using supervised bilevel learning, implicit upper-level
gradients, and training velocity models. Our aim is complementary: we study an outer design loop in which the design objective itself is
interchangeable. We compare an oracle design objective, which requires the true velocity model, with practical and less expensive alternatives.


Our application is motivated by guided-wave thickness mapping and scanning laser
Doppler vibrometry (LDV). Previous work has shown that FWI can reconstruct plate thickness from guided-wave measurements~\citep{zuo2022quantitative}, making it relevant to non-destructive testing of thin structures such as plates. LDV enables
highly flexible receiver placement, but a measurement at each spatial position must be repeated many times and stacked to improve its signal-to-noise ratio~\citep{Thomsen2024}.
Acquisition time therefore grows quickly with the number of positions, strongly motivating OED for a reduced acquisition
layout.

This motivates an outer design loop for online acquisition rather than a
single layout designed before the experiment. In practice, the LDV
first measures the wavefield at the current receiver coordinates. A short block of FWI is then performed. The selected design objective uses that partially
updated velocity model to calculate new receiver coordinates. The LDV then moves to
those coordinates, acquires a new stacked data set, and the loop repeats. All
design objectives studied below are embedded in this same outer design loop.
Figure~\ref{fig:diff-oed-workflow} illustrates its
measurement--inversion--design--reacquisition sequence.

% This study presents a proof of concept with three main contributions:
% \begin{itemize}
%   \item a fully differentiable acquisition and acoustic FWI workflow that optimizes continuous receiver coordinates;
%   \item a comparison of an oracle model error design objective with a waveform misfit objective and a stochastic illumination objective in the same computational framework; and
%   \item preliminary evidence that the compared objectives improve reconstruction in a synthetic non-destructive testing setting.
% \end{itemize}

Our contributions are a differentiable acoustic FWI workflow for continuous
receiver placement, a comparison of oracle model error, waveform misfit, and stochastic
illumination objectives, and a synthetic demonstration that each improves
reconstruction.

% More broadly, differentiating the complete acquisition and inversion workflow
% makes its components modular: expensive solver or reconstruction components
% could in principle be replaced by learned surrogates, while internal
% algorithmic hyperparameters could be optimized jointly with the acquisition geometry. This study does not yet explore these possibilities.

\section{Differentiable acquisition design}

\begin{figure}[t]
  \centering
  \includegraphics[width=\linewidth]{figures/differentiable-oed-double-loop-draft_ipnut_output.pdf}
  \caption{Nested differentiable OED. The inputs to the outer design loop are the initial receiver positions and initial velocity model. At outer design iteration $t$, measurements at
$\vect{p}^t$ drive $K$ FWI updates; the design objective then updates the
receivers to $\vect{p}^{t+1}$, where data are reacquired. The outputs of this workflow are the updated receiver positions and the corresponding reconstructed models. Red crosses denote fixed sources and yellow stars denote movable
  receivers.}
  \label{fig:diff-oed-workflow}
\end{figure}

\subsection{Full-waveform inversion}

Figure~\ref{fig:diff-oed-workflow} shows the bilevel structure of the workflow. The two optimization levels are indexed separately: $t$
indexes iterations of the outer design loop, and $k$ indexes FWI updates. Let
$\vect{m}$ denote the discretized velocity model,
$\vect{p}\in\mathbb{R}^{2N_r}$ the continuous receiver coordinates of $N_r$ receivers,
and $\mathcal{F}_{\vect{p}}(\vect{m})$ the time-domain acoustic forward
operator sampled at those coordinates. Let $\mathcal{U}_{m}$ denote the FWI
optimizer. For predicted data
$\vect{d}^{\mathrm{pred}}=\mathcal{F}_{\vect{p}}(\vect{m})$, the data misfit is
the mean squared error between observed and predicted waveforms,
\begin{equation}\label{eq:fwi_loss}
  \Ldata(\vect{d}^{\mathrm{pred}},\vect{d}^{\mathrm{obs}})
  = \left\|\vect{d}^{\mathrm{pred}}-\vect{d}^{\mathrm{obs}}\right\|_2^2.
\end{equation}

At outer design iteration $t$, the geometry is
$\vect{p}^{t}$ and the corresponding observed waveforms are
$\vect{d}^{\mathrm{obs}}(\vect{p}^{t})$. Starting from $\vect{m}_0^{t}$, $K$
steps of FWI produce
\begin{align}
  \vect{m}_{k+1}^{t}
  &= \mathcal{U}_{m}\!\left(
     \vect{m}_k^{t},
     \nabla_{\vect{m}}\Ldata\!\left(
       \mathcal{F}_{\vect{p}^{t}}\!\left(
         \vect{m}_k^{t}
       \right),
       \vect{d}^{\mathrm{obs}}(\vect{p}^{t})
     \right)
  \right),
  \qquad k=0,\ldots,K-1,
  \label{eq:fwi-update}
\end{align}
The latest velocity model is carried over to the next outer design iteration,
such that $\vect{m}_0^{t+1}=\vect{m}_K^{t}$. The complete workflow is differentiable and uses automatic differentiation to compute the FWI gradients as well as the outer design gradients.

\subsection{Interchangeable design objectives}

Here, we present the three design objectives used in this study. Each objective
uses the current receiver coordinates $\vect{p}^{t}$ and the corresponding
reconstruction state defined below. The objectives are interchangeable, and the
same workflow can be used to compare their performance.

\paragraph{Oracle model error.}
For synthetic validation, the privileged objective directly measures
reconstruction quality against the true model $\vect{m}^{\star}$,
\begin{equation}
  \mathcal{L}_{\mathrm{model}}(\vect{p}^{t})
  = \left\|\vect{m}_K^{t}-\vect{m}^{\star}\right\|_2^2.
\end{equation}
It provides an upper benchmark but is unavailable in a physical experiment.

\paragraph{Waveform misfit.}
Let $\vect{d}^{\mathrm{obs}}(\vect{p})$ denote measurements acquired after the
receivers have moved to $\vect{p}$. We evaluate the reconstructed model using
the same acquisition geometry:
\begin{equation}
  \mathcal{L}_{\mathrm{wave}}(\vect{p}^{t})
  = \left\|\mathcal{F}_{\vect{p}^{t}}
  \!\left(\vect{m}_K^{t}\right)
  -\vect{d}^{\mathrm{obs}}(\vect{p}^{t})\right\|_2^2.
  \label{eq:waveform}
\end{equation}
This objective does not require the unknown true material model. Because the
waveform misfit is evaluated at a different acquisition geometry at every
outer design iteration, its values across design iterations are not expected to
decrease monotonically.

\paragraph{Stochastic illumination.}
The illumination objective scores a receiver geometry by how strongly and how
uniformly the simulated data respond to perturbations throughout the model.
It uses the diagonal of the local Gauss--Newton matrix $J^\top J$, where
$J=\partial\mathcal{F}_{\vect{p}}(\vect{m})/\partial\vect{m}$ is evaluated at
the current reconstructed model. Large diagonal entries indicate model
locations to which the acquisition is locally sensitive. We estimate this
diagonal with random probes and vector--Jacobian products, so the full Jacobian
and Gauss--Newton matrix are never formed. The loss penalizes spatial
nonuniformity while rewarding a large mean sensitivity. This criterion requires neither the true model nor observed
waveforms. Appendix~\ref{app:illumination-derivation} derives the estimator and
gives the exact scalarization; its numerical settings are listed in
Table~\ref{tab:outer-design-parameters}.

\subsection{Design update and differentiation paths}

Let $\Ldesign$ denote one of the three objectives above and $\mathcal{U}_{p}$
the coordinate optimizer. The receiver-coordinate update is
\begin{equation}
  \vect{p}^{t+1}
  = \mathcal{U}_{p}\!\left(
       \vect{p}^{t},
       \nabla_{\vect{p}}\Ldesign(\vect{p}^{t})
     \right),
  \label{eq:design-update}
\end{equation}
The objectives induce two different differentiation paths. For the oracle
model error and waveform misfit
objectives, reverse-mode automatic differentiation propagates the
receiver-coordinate gradient through the complete unrolled FWI algorithm,
including all $K$ FWI updates and their wave simulations. By contrast, the
illumination gradient is differentiated through the forward wave solver and
its vector--Jacobian products, but not through the FWI trajectory.  Illumination therefore reduces the additional graph required to
calculate the design update. All methods
still perform $K$ FWI updates to advance the reconstruction between design
iterations. We implement the workflow in JAX using a differentiable wave solver
in the spirit of j-Wave~\citep{stanziola2023jwave}.

We do not assume that the inner FWI problem is fully converged. Instead, the
$K$ updates in Eq.~\eqref{eq:fwi-update} define the reconstruction presented to
the outer objective, so the design is optimized for a given number of FWI
updates. We fix $K$ in this study. In future, the inversion budget $K$, or a
learned stopping rule based on the level of convergence, could itself be
optimized. In this finite-iteration formulation, the chosen optimizer and
stopping point are hyperparameters, and retaining the updates determines the
memory cost of differentiation.

\section{Experimental setup}

\paragraph{Model and wave simulation.}
Following the acoustic proof of concept described above, the observations are generated and FWI is performed with the same differentiable two-dimensional 
acoustic solver. We use a $25\,\mathrm{kHz}$ source in a homogeneous background
containing one faster square inclusion. The synthetic inclusion also has a
density contrast, whereas FWI reconstructs wave speed only from a homogeneous
initial model and holds density fixed. This choice reduces the computational cost of FWI. Wave propagation uses a first-order
staggered-grid scheme with bilinear source injection and receiver sampling.
Appendix Table~\ref{tab:acoustic-parameters} gives the complete discretization,
material, boundary, and preprocessing parameters.

\paragraph{Acquisition geometry.}
Four sources remain fixed around the target, while the coordinates of 30 receivers are continuously optimized. The receivers begin on a diagonal line crossing the
model and are projected to the non-absorbing region after every outer
design iteration. Appendix Fig.~\ref{fig:true-model-geometry} shows the true
model and nominal initial geometry; exact coordinates and bounds are listed in
Appendix Table~\ref{tab:acquisition-parameters}.

\paragraph{Inversion and design optimization.}
% As defined in Eq.~\eqref{eq:fwi_loss}, FWI minimizes the mean squared difference between trace-normalized observed and predicted waveforms. 

For all comparisons, we use $T=5$ outer design iterations with $K=20$ FWI updates per design iteration. Each
optimized reconstruction therefore uses $TK=100$ FWI updates, while the
fixed-layout baseline receives one uninterrupted 100-update inversion. The
methods thus have the same FWI reconstruction budget, although their design costs
differ.

The FWI algorithm uses stochastic gradient descent (SGD). Compared with Adam,
SGD has lower memory requirements, making it practical to retain the computation
graph for the 20 FWI updates required by the oracle and waveform misfit objectives. Appendix Tables~\ref{tab:fwi-parameters}
and~\ref{tab:outer-design-parameters} report the remaining optimizer,
illumination, and randomization settings.

% We compare the fixed initial layout with layouts optimized using the oracle
% model-error, waveform misfit, and illumination objectives. 

% The primary metric is the mean squared model error over the complete grid. The
% multi-seed protocol holds the target fixed while varying the paired receiver
% initialization and source schedule; wall-clock time and peak memory are recorded
% for each run.

\section{Preliminary results}

We compare the fixed layout with layouts optimized using the oracle model
error, waveform misfit, and illumination objectives. The primary metric is the mean squared model error (MSE) over the complete grid. The multi-seed protocol holds the true and initial models fixed while varying the paired receiver initialization and stochastic FWI source sequence.

% Across ten paired repetitions, every optimized layout reduces final model MSE
% relative to its fixed-layout baseline (Table~\ref{tab:multiseed-summary}). The
% oracle and illumination objectives give mean gains of $18.6\%$ and $16.7\%$,
% respectively, with illumination showing greater variability. Illumination also has the lowest runtime and memory usage, with a mean runtime of \(356.5\,\mathrm{s}\) and a mean peak memory increase of \(9.8\,\mathrm{GiB}\), compared with \(479.2\,\mathrm{s}\) and \(64.3\,\mathrm{GiB}\) for oracle model error and \(503.1\,\mathrm{s}\) and \(66.1\,\mathrm{GiB}\) for waveform misfit. Waveform misfit
% gives a smaller mean gain of $9.4\%$ but the lowest variability. Individual repetitions are
% shown in Appendix Fig.~\ref{fig:paired-improvement}. An illustrative
% reconstruction is shown in Fig.~\ref{fig:reconstruction-comparison}, and the
% corresponding receiver trajectories appear in Appendix
% Fig.~\ref{fig:receiver-trajectories}. The reconstruction using the oracle model error is also shown in Appendix Fig.~\ref{fig:reconstruction-comparison-oracle}.

Across ten paired repetitions, every optimized layout reduces final model MSE
relative to its fixed-layout baseline (Table~\ref{tab:multiseed-summary}).
Oracle and illumination give the largest mean improvements, while illumination
requires substantially less runtime and memory than the objectives
differentiated through FWI. Waveform misfit gives a smaller
but less variable improvement. Figure~\ref{fig:reconstruction-comparison} shows
one representative reconstruction; individual repetitions and receiver
trajectories are reported in the Appendix.


\begin{table}[!htbp]
  \centering
  \scriptsize
  \setlength{\tabcolsep}{3.5pt}
  \caption{Multi-seed reconstruction and runtime summary. Values are
  mean $\pm$ sample standard deviation; improvements are relative to the paired
  fixed-layout run. Peak memory increase is reported as a mean rounded to
  $0.1\,\mathrm{GiB}$. Every reconstruction uses 100 FWI updates.}
  \label{tab:multiseed-summary}
  \begin{tabular}{@{}lrrrrr@{}}
    \toprule
    Method & Final model MSE & Improvement (\%) & Improved & Runtime (s) & Peak $\Delta$Memory (GiB) \\
    \midrule
    Fixed layout & $3358.9\pm65.5$ & $0.0$ & - & $166.3\pm0.9$ & 1.2 \\
    Stochastic illumination & $2796.2\pm274.6$ & $16.7\pm8.6$ & 10/10 & $356.5\pm0.8$ & 9.8 \\
    Waveform misfit & $3041.6\pm125.4$ & $9.4\pm4.3$ & 10/10 & $503.1\pm0.7$ & 66.1 \\
    Oracle model error & $2732.3\pm242.4$ & $18.6\pm7.1$ & 10/10 & $479.2\pm1.2$ & 64.3 \\
    \bottomrule
  \end{tabular}
\end{table}

\begin{figure}[t]
  \centering
  \includegraphics[width=0.85\linewidth]{figures/comparison_plot_full_auto_diff_cropped_colorbar_true_model_outlined.pdf}
  \caption{Reconstruction and receiver-geometry comparison for one
  representative repetition. The dashed cyan box marks the true
  $55\,\mathrm{mm}$ square inclusion; red crosses are the four fixed sources and
  yellow stars are the 30 receivers. The wave-speed scale is shared, and
  percentages give MSE reduction from the fixed layout. Every reconstruction
  uses 100 FWI updates. The color bar is cropped at 80\% of the maximum value of the true model to emphasize differences between the design methods.}
  \label{fig:reconstruction-comparison}
\end{figure}

\section{Discussion and limitations}

The experiments show that differentiable OED can improve continuous receiver
placement for a finite nonlinear inversion. More broadly, the approach addresses
one source of misspecification in practical FWI: the acquisition geometry must
usually be chosen before the informative measurement locations are known. By
alternating inversion, receiver updates, and reacquisition, the proposed loop
can adapt an initially suboptimal geometry as information about the model
emerges.

The three design objectives should be viewed as different compromises rather than a
definitive ranking. The oracle model error design objective is a privileged reference that is unavailable in a physical experiment. Illumination
currently approaches the oracle’s mean reconstruction benefit at lower computational
cost, but with greater variation across repetitions. Waveform misfit is less
variable but gives a smaller mean improvement. The waveform misfit objective assumes that receiver layouts producing a better data fit after a finite FWI block also produce velocity models closer to the true model. This assumption is not guaranteed in a nonlinear, ill-posed inverse problem.

The present conclusions are limited to noise-free synthetic observations
generated with the same solver used for inversion. Robustness to measurement
noise and modeling error therefore remains untested. Moreover, the methods are
matched by their number of FWI updates rather than by total computational or
physical acquisition cost, and the learned layouts depend on the chosen
optimizers and update budgets. Future work should evaluate independent forward
models, LDV-calibrated noise, fixed acquisition budgets, and ultimately
closed-loop laboratory measurements.

% Illumination avoids differentiation through the unrolled FWI trajectory,
% substantially reducing runtime and peak memory. Its scaling to longer
% trajectories remains untested, and larger problems may require checkpointing,
% truncated unrolling, or implicit gradients.

% The methods are matched by 100 FWI updates rather than total computation or
% physical acquisition time. Future LDV studies should instead compare
% reconstruction quality under fixed acquisition and computational budgets.

% The observations are noise-free and generated with the same solver used for
% inversion, so robustness to noise and modeling error remains untested. Future
% evaluations should use LDV-calibrated noise, an independent forward model, and
% ultimately closed-loop laboratory acquisition.

Overall, the results support differentiable acquisition as a flexible way to
decide not only how to reconstruct a model, but also where the data needed for
that reconstruction should be measured. The present approach provides also an exciting route for
investigating which hyperparameters, optimizers, and wave solvers could
themselves become optimization variables.





% The optimized layouts also depend on the chosen optimizer and update budget;
% these choices could themselves become optimization variables.

\clearpage
\bibliographystyle{plainnat}
\bibliography{references}

% Sim2Science requires the checklist after the references and before the
% appendix. It does not count toward the five-page limit.
\newpage
\input{checklist.tex}

\clearpage
\appendix
\makeatletter
\setlength{\@fptop}{0pt}
\makeatother
\section{Stochastic illumination criterion}
\label{app:illumination-derivation}

This section derives the stochastic estimator used by the illumination
objective. For source $s$, let
\begin{equation}
  J_s^t
  = \left.
    \frac{\partial \mathcal{F}_{\vect{p}}^{(s)}(\vect{m})}
         {\partial \vect{m}}
    \right|_{\vect{p}=\vect{p}^t,\,\vect{m}=\vect{m}_0^t}
\end{equation}
be the forward-model Jacobian at outer design iteration $t$. Its local
Gauss--Newton matrix is $(J_s^t)^\top J_s^t$, whose diagonal measures the
squared data sensitivity of each model parameter.

Let $\vect{z}_{s,j}\in\{-1,+1\}^{D_s}$ be an independent Rademacher vector,
where $D_s$ is the number of data samples for source $s$. Since
$\mathbb{E}[\vect{z}_{s,j}\vect{z}_{s,j}^\top]=I$, we have
\begin{align}
  \mathbb{E}\!\left[
    ((J_s^t)^\top\vect{z}_{s,j})
    \odot
    ((J_s^t)^\top\vect{z}_{s,j})
  \right]
  &= \operatorname{diag}\!\left(
       (J_s^t)^\top
       \mathbb{E}[\vect{z}_{s,j}\vect{z}_{s,j}^\top]
       J_s^t
     \right) \\
  &= \operatorname{diag}\!\left((J_s^t)^\top J_s^t\right),
  \label{eq:illumination-expectation}
\end{align}
where $\odot$ denotes elementwise multiplication. This identity gives an
unbiased randomized diagonal estimator~\citep{bekas2007diagonal}.

In the implementation, each probe is normalized as
$\vect{q}_{s,j}=\vect{z}_{s,j}/\sqrt{D_s}$. With $S$ sources and $Q$ probes per
source, the unsmoothed illumination estimate is
\begin{equation}
  \widehat{\vect{h}}
  = \frac{1}{SQ}\sum_{s=1}^{S}\sum_{j=1}^{Q}
    ((J_s^t)^\top\vect{q}_{s,j})
    \odot
    ((J_s^t)^\top\vect{q}_{s,j}).
  \label{eq:illumination}
\end{equation}
Thus, $\mathbb{E}[\widehat{\vect{h}}]$ is the source-averaged Gauss--Newton
diagonal scaled by $1/D_s$ (the data size is identical for all sources here).
Each product $(J_s^t)^\top\vect{q}_{s,j}$ is computed as a
vector--Jacobian product by automatic differentiation, avoiding explicit
construction of $J_s^t$.

After one spatial smoothing application, denote the resulting field by
$\vect{h}$ and its spatial mean by $\bar h$. The receiver coordinates minimize
\begin{equation}
  \mathcal{L}_{\mathrm{illum}}(\vect{p}^t;\vect{m}_0^t)
  = \operatorname{Var}\!\left(\frac{\vect{h}}{\bar h+\epsilon}\right)
    - \beta\log(\bar h+\epsilon),
  \qquad \beta=0.05,\quad \epsilon=10^{-12}.
  \label{eq:illumination-scalarization}
\end{equation}
The first term penalizes relative spatial variation and therefore favors even
coverage; the second decreases as the mean illumination grows and therefore
prevents a uniformly weak geometry from being preferred. The probes are fixed
across design iterations within a repetition, making the optimized stochastic
objective deterministic within that run.

\section{Experimental parameters}

The main text summarizes the experimental logic. Tables~\ref{tab:acoustic-parameters}--
\ref{tab:outer-design-parameters} collect the numerical values used for the
reported experiments.

\begin{figure}[!htbp]
  \centering
  \includegraphics[width=0.68\linewidth]{figures/sh0_model_sources_receivers_outlined.pdf}
  \caption{True wave-speed model and nominal initial acquisition geometry.
  The white square is the $55\,\mathrm{mm}$ inclusion, red crosses are the four
  fixed sources, and yellow stars are the 30 initial receiver positions.}
  \label{fig:true-model-geometry}
\end{figure}

\begin{table}[!htbp]
  \centering
  \small
  \caption{Acoustic model and wave-simulation parameters.}
  \label{tab:acoustic-parameters}
  \begin{tabular}{@{}L{0.34\linewidth}L{0.59\linewidth}@{}}
    \toprule
    Parameter & Value \\
    \midrule
    Governing model & Two-dimensional scalar acoustic equation; first-order staggered-grid time stepping \\
    Domain and grid & $1.1\,\mathrm{m}\times1.1\,\mathrm{m}$; $210\times210$ points; $\Delta x=\Delta y=5.26\,\mathrm{mm}$ \\
    Background material & Wave speed $2100\,\mathrm{m\,s^{-1}}$; density $1180\,\mathrm{kg\,m^{-3}}$ \\
    True inclusion & $55\,\mathrm{mm}$ square centered at $(-50,50)\,\mathrm{mm}$; wave speed $3200\,\mathrm{m\,s^{-1}}$; density $3000\,\mathrm{kg\,m^{-3}}$ \\
    Inverted parameters & Wave speed only; homogeneous initial speed and homogeneous background density \\
    Source wavelet & Ricker pulse; central frequency $25\,\mathrm{kHz}$; maximum modeled frequency $50\,\mathrm{kHz}$ \\
    Spatial resolution & Eight grid points per minimum wavelength at $50\,\mathrm{kHz}$ in the background \\
    Time sampling & 807 steps; $\Delta t=0.930\,\mu\mathrm{s}$; record length $750.8\,\mu\mathrm{s}$ \\
    Absorbing boundary & 19-cell ($0.1\,\mathrm{m}$) convolutional PML; target reflection $10^{-6}$; polynomial order 3; maximum stretching coefficient 3 \\
    Source/receiver sampling & Bilinear injection and sampling \\
    Trace preprocessing & Each receiver trace divided by its maximum absolute amplitude, with floor $10^{-12}$ \\
    \bottomrule
  \end{tabular}
\end{table}

\begin{table}[!htbp]
  \centering
  \small
  \caption{Acquisition-geometry parameters.}
  \label{tab:acquisition-parameters}
  \begin{tabular}{@{}L{0.34\linewidth}L{0.59\linewidth}@{}}
    \toprule
    Parameter & Value \\
    \midrule
    Sources & Four fixed sources at $(\pm0.35,0)\,\mathrm{m}$ and $(0,\pm0.35)\,\mathrm{m}$ \\
    Receivers & 30 receivers with two continuously optimized coordinates each \\
    Initial receiver geometry & Equally spaced on the diagonal from $(-0.15,-0.15)\,\mathrm{m}$ to $(0.15,0.15)\,\mathrm{m}$, with repetition-specific jitter in the multi-seed study \\
    Admissible receiver region & Non-absorbing region $[-0.45,0.45]^2\,\mathrm{m}$; projection during objective evaluation and after outer design iterations \\
    \bottomrule
  \end{tabular}
\end{table}

\begin{table}[!htbp]
  \centering
  \small
  \caption{Inner FWI parameters.}
  \label{tab:fwi-parameters}
  \begin{tabular}{@{}L{0.34\linewidth}L{0.59\linewidth}@{}}
    \toprule
    Parameter & Value \\
    \midrule
    Data loss & Mean squared error between trace-normalized observed and predicted waveforms \\
    Model optimizer & Stochastic gradient descent; notebook learning-rate parameter 25, multiplied by $10^{6}$ before optimizer construction \\
    Stochastic source sampling & One of four sources drawn uniformly per FWI update; repetition-specific schedule shared across methods \\
    Model-gradient processing & One spatial smoothing application per FWI update, followed by multiplication by 10 \\
    Optimized-method budget & $T=5$ outer design iterations; $K=20$ FWI updates per iteration; $TK=100$ total FWI updates \\
    Fixed-geometry budget & One uninterrupted 100-update FWI inversion \\
    Model carry & $\vect{m}^{t+1}_0=\vect{m}^{t}_K$ between outer design iterations \\
    \bottomrule
  \end{tabular}
\end{table}

\begin{table}[!htbp]
  \centering
  \small
  \caption{Outer design loop and illumination parameters.}
  \label{tab:outer-design-parameters}
  \begin{tabular}{@{}L{0.34\linewidth}L{0.59\linewidth}@{}}
    \toprule
    Parameter & Value \\
    \midrule
    Outer design gradient & Reverse-mode automatic differentiation; reconstruction-aware objectives differentiate through the 20 unrolled SGD updates and receiver sampling \\
    Receiver optimizer & Optax Adam; learning rate 4; default $(\beta_1,\beta_2)=(0.9,0.999)$ \\
    Coordinate constraints & Project candidate and updated coordinates to the admissible region \\
    Illumination gradient & Automatic differentiation through the wave-solver VJPs and illumination scalarization; no unrolled FWI derivative \\
    Illumination probes & 12 Rademacher probes per source, each divided by the square root of its number of entries; repetition-specific probes fixed within each run \\
    Illumination processing & Average over probes and sources; one spatial smoothing application; normalize by mean illumination \\
    Illumination scalarization & Uniformity variance plus $0.05$ times negative log mean strength; numerical offset $10^{-12}$ \\
    \bottomrule
  \end{tabular}
\end{table}

\clearpage
\section{Multi-seed and computational results}

Each repetition uses independent random seeds controlling three sources of
variation:
\begin{itemize}
  \item \textbf{Layout:} the perturbation of the initial receiver coordinates;
  \item \textbf{FWI sources:} the stochastic sequence of sources selected by
  the inner FWI updates; and
  \item \textbf{Illumination probes:} the stochastic probes used by the
  illumination objective.
\end{itemize}
Within each paired repetition, all methods share the same initial receiver
layout and FWI source sequence, while the illumination-probe seed applies only
to the illumination method.

\paragraph{Computational measurements.}
The runs were executed under Linux on an x86-64 workstation equipped with one
AMD Ryzen Threadripper 3970X processor (32 physical cores, 64 hardware threads,
and a reported frequency range of $2.2$--$3.7\,\mathrm{GHz}$). Total runtime is
the wall-clock time for a complete method run, including 100 FWI updates and,
for the optimized methods, five design updates. For the oracle and waveform
misfit objectives, the nested FWI and receiver-hypergradient calculations are
compiled and timed together and cannot be separated; illumination records its
design and subsequent FWI timings separately. The saved measurements do not
separate JAX compilation or warm-up from execution, so the reported runtimes
are end-to-end measurements under the compilation and cache state of each run.
The peak (memory) RSS increase is the maximum recorded process resident set size minus
the process RSS at the start of the run.

\begin{figure}[!htbp]
  \centering
  \includegraphics[width=0.90\linewidth]{figures/paired_improvement_outlined.pdf}
  \caption{Improvement in final model MSE relative to the paired fixed-layout
  baseline. Points are individual completed repetitions, diamonds show the
  mean, and bars show one sample standard deviation. The dashed line denotes no
  improvement. Each method is evaluated over ten paired repetitions.}
  \label{fig:paired-improvement}
\end{figure}

% \begin{figure}[!htbp]
%   \centering
%   \includegraphics[width=0.92\linewidth]{figures/paired_model_mse.pdf}
%   \caption{Final model MSE for the fixed layout and optimized objectives. Gray
%   lines connect results belonging to the same repetition; lower values are
%   better.}
%   \label{fig:paired-model-mse}
% \end{figure}

\begin{figure}[!htbp]
  \centering
  \includegraphics[width=0.84\linewidth]{figures/runtime_vs_improvement_outlined.pdf}
  \caption{Total runtime and paired reconstruction improvement for every
  optimized run. The illumination objective is faster but more variable than
  the objectives differentiated through the unrolled FWI trajectory.}
  \label{fig:runtime-improvement}
\end{figure}

\begin{figure}[!htbp]
  \centering
  \includegraphics[width=\linewidth]{figures/compute_time_and_memory_per_run_outlined.pdf}
  \caption{Total runtime and additional peak host memory by method. Points are
  individual runs, diamonds show the mean, and bars show one sample standard
  deviation. Differentiating through the unrolled inversion substantially
  increases memory relative to the illumination differentiation path.}
  \label{fig:compute-memory}
\end{figure}

\clearpage
\section{Oracle model error reconstruction}

Figure~\ref{fig:reconstruction-comparison-oracle} shows the same reconstruction comparison as Fig.~\ref{fig:reconstruction-comparison} but for the oracle model error objective.

\begin{figure}[ht]
  \centering
  \includegraphics[width=0.85\linewidth]{figures/comparison_plot_full_auto_diff_cropped_colorbar_outlined.pdf}
  \caption{Reconstruction and receiver-geometry comparison for one
  representative repetition. The dashed cyan box marks the true
  $55\,\mathrm{mm}$ square inclusion; red crosses are the four fixed sources and
  yellow stars are the 30 receivers. The wave-speed scale is shared, and
  percentages give MSE reduction from the fixed layout. Every reconstruction
  uses 100 FWI updates. The color bar is cropped at 80\% of the maximum value of the true model to emphasize differences between the design methods.}
  \label{fig:reconstruction-comparison-oracle}
\end{figure}

\clearpage
\section{Receiver trajectories}

Figure~\ref{fig:receiver-trajectories} shows how the receiver geometries evolve
over five outer design iterations for each optimized objective.

\begin{figure}[ht]
  \centering
  \includegraphics[width=0.96\linewidth]{figures/comparison_plot_full_auto_diff_trajectories_outlined.pdf}
  \caption{Receiver motion during acquisition optimization. Panel (a) shows
  the fixed initial layout. Panels (b)--(d) show the layouts optimized with the
  illumination, waveform misfit, and oracle model error
  objectives, respectively. Orange polylines connect each receiver's successive
  positions over the five outer design iterations; yellow stars mark the final receiver
  positions. Red crosses are the fixed sources, and the dashed cyan square marks
  the true inclusion.}
  \label{fig:receiver-trajectories}
\end{figure}

\end{document}
